% Generated draft: requires mathematical and source review.
\documentclass[11pt]{article}
\usepackage[margin=25mm]{geometry}
\usepackage{fontspec}
\setmainfont{Latin Modern Roman}
\usepackage{amsmath,amssymb,mathtools,unicode-math}
\setmathfont{Latin Modern Math}
\usepackage{enumitem,needspace,xurl,hyperref}
\setlength{\parindent}{0pt}
\setlength{\parskip}{5pt}
\setlength{\emergencystretch}{3em}
\newcommand{\N}{\mathbb{N}}
\newcommand{\Z}{\mathbb{Z}}
\newcommand{\Q}{\mathbb{Q}}
\newcommand{\R}{\mathbb{R}}
\newcommand{\C}{\mathbb{C}}
\newcommand{\F}{\mathbb{F}}
\newcommand{\A}{\mathbb{A}}
\newcommand{\PP}{\mathbb{P}}
\newcommand{\E}{\mathbb{E}}
\newcommand{\eps}{\varepsilon}
\newcommand{\dd}{\,\mathrm{d}}
\begin{document}
\section*{UnsolvedMath ID: 30003704}
\textbf{Source problem code:} catalogue entry
\subsection{Definitions and mathematical statement}
The supplied question concerns finite generation of graph-configuration homology as a functor. Its precise formulation requires a category \(\mathcal C\) of graphs, a coefficient ring \(R\), a particle number \(k\), a convention for configuration spaces, and compatible maps on homology. These data are not specified in the supplied record.

For a topological realization of a graph \(G\) and a nonnegative integer \(k\), the usual ordered configuration space is
\[
\{(x_1,\ldots,x_k)\in G^k : x_i\ne x_j\text{ whenever }i\ne j\}.
\]
The usual unordered configuration space is its quotient by the symmetric group permuting coordinates. The notation \(\operatorname{Conf}_k(G)\) in the source does not specify which convention is intended. The allowed graphs and graph morphisms must also be specified so that the displayed assignment actually defines a functor.

If the intended meaning is finite generation of a functor \(F:\mathcal C\to R\text{-}\mathrm{Mod}\), its usual module interpretation is the existence of finitely many objects \(G_1,\ldots,G_m\) and elements \(v_i\in F(G_i)\) such that, for every object \(G\),
\[
F(G)=\operatorname{span}_R\{F(f)(v_i):1\le i\le m,\ f\in\operatorname{Hom}_{\mathcal C}(G_i,G)\}.
\]
This definition is conditional: the record does not establish that this is the precise finite-generation notion used in the cited result.

Subject to recovering the hypotheses and conventions of that result, the supplied extension question is: for the same fixed particle number \(k\), graph category, coefficients, and functorial structure, is
\[
F_{q,k}:G\longmapsto H_q(\operatorname{Conf}_k(G);R)
\]
finitely generated for every integer \(q\ge0\)? Here \(R\) records an unresolved coefficient choice, rather than an additional assertion about arbitrary coefficients. The record reports an existing finite-generation result in degrees \(q=0,1\) and asks whether it extends to all nonnegative degrees. It does not specify the quantification over \(k\), or whether any uniformity in \(k\) or \(q\) is intended.

The supplied dated assessment also reports finite generation for certain FI-graph families, but does not identify those families or their hypotheses. Those qualifications must be recovered before a fully determined statement can be made.
\subsection{Short English statement}
The question asks whether the reported finite-generation theorem for graph-configuration homology in degrees zero and one extends to every nonnegative homological degree, with the original hypotheses retained. The graph category, morphisms, coefficients, configuration-space convention, particle-number quantifiers, and relevant FI-graph families are missing, so this is a conditional statement requiring mathematical review.
\subsection{Sources}
\begin{itemize}
\item Supplied UnsolvedMath local registry record, permanent ID \(30003704\), source problem code \texttt{catalogue entry}, entitled \emph{Finite Generation of Graph-Configuration Homology Functors}. The record lists proposed year 2018 and identifies no individual proposer.
\item Supplied report citation: \emph{Topology of Arrangements and Representation Stability (2019)}, Oberwolfach Reports, DOI \texttt{10.4171/owr/2018/2}, pp.\ 43--123. \url{https://doi.org/10.4171/owr/2018/2}. The background describes Oberwolfach Reports as documenting mathematical workshops and their open-problem discussions; no exact passage or theorem specifying the omitted hypotheses is supplied.
\item Supplied literature reference: Ramos, \emph{FI-algebras and graph configuration spaces}. \url{https://arxiv.org/abs/1805.05316}.
\item Literature assessment dated exactly \texttt{2026-08-21}: the supplied curation labels the problem open and reports finite generation in homological degrees at most one and for the stated FI-graph families, while reporting the general higher-degree question as open. This is supplied dated curation, not an independently verified literature assessment. Its time-sensitive status and cited mathematical claims require independent verification.
\end{itemize}
\end{document}
